% Concise theory for both papers; full derivations are preserved in
% paper/audits/streamlining_20260911/full_theory.tex.
\section{Mode Collapse and Verified Replay}
\label{app:theory}

This appendix establishes the categorical claims behind the method: centered
binary-reward estimators amplify a pre-existing imbalance among correct modes,
whereas persistent verified replay supplies a loss barrier for banked modes.
The complete-response argument then states what would be needed to transfer
that protection to a language model. These results concern idealized update
rules; they do not certify the implemented AdamW/PPO trajectories.

\subsection{Categorical reduction and mean updates}
\label{app:theory-grpo-mean}

Fix one prompt and a finite response alphabet
$\mathcal A=\mathcal C\mathbin{\dot\cup}\mathcal N$, with $m\ge2$ correct
execution modes and $n\ge1$ incorrect categories. Let
$p_a=\operatorname{softmax}(z)_a$, $P=\sum_{c\in\mathcal C}p_c$, and
$q_c=p_c/P$. Thus $P$ measures correctness and $q$ the distribution among
correct modes. Throughout the mean-flow analysis, logits are finite initially,
$0<P(0)<1$, explicit reference KL is zero, and the prompt is isolated.

Draw a group $A_1,\ldots,A_G\stackrel{\mathrm{iid}}{\sim}p$ with $G\ge2$,
write $R_i=\1[A_i\in\mathcal C]$, and let $M=\sum_iR_i$. Both GRPO variants
can be written, at the on-policy point where the PPO ratio equals one, as
\begin{equation}
  \widehat g_G
  =\frac1G\sum_{i=1}^G
   w(M)\left(R_i-\frac{M}{G}\right)\nabla_z\log p_{A_i}.
  \label{eq:group-estimator}
\end{equation}
For Dr.GRPO, $w(M)=1$ after absorbing the shared positive length scale. For
binary-reward GRPO, $w(M)$ is the reciprocal group reward standard deviation
on mixed groups. We absorb the same fixed response-length normalizer in both
variants; response-dependent row normalization is outside this reduction.
In either case $0<w(\ell)<\infty$ for $\ell\in\{1,\ldots,G-1\}$.
We define the entire update as zero on all-equal groups, without evaluating an
undefined reciprocal standard deviation; likewise, $w(M)M(G-M)$ below is
defined as zero when $M=0$ or $M=G$. We take explicit reference KL to be zero, as
in the experiments.

\begin{lemma}[A group baseline does not remove frequency amplification]
\label{lem:grpo-mean}
For $0<P<1$, the expected group update is
\begin{equation}
  \E[\widehat g_G]
  =c_G(P)\nabla_z P,
  \qquad
  c_G(P)=
  \frac{\E\!\big[w(M)M(G-M)\big]}{G^2P(1-P)}>0.
  \label{eq:expected-group-gradient}
\end{equation}
For Dr.GRPO, $c_G(P)=(G-1)/G$. Hence reward centering and binary reward
standardization change the speed, but not the expected direction, of the
binary-reward policy gradient.
\end{lemma}

\begin{proof}
Condition on the complete binary reward vector, whose sum is $M$. The score
identities give
\[
  \E[\nabla_z\log p_{A_i}\mid R_i=1]=\nabla_z\log P,
  \qquad
  \E[\nabla_z\log p_{A_i}\mid R_i=0]
  =\nabla_z\log(1-P).
\]
There are $M$ correct rows with centered reward $(G-M)/G$ and $G-M$
incorrect rows with centered reward $-M/G$. Therefore the conditional
expectation of the unaveraged sum in~\eqref{eq:group-estimator} is
\[
  w(M)\frac{M(G-M)}{G}
  \left(\nabla_z\log P-\nabla_z\log(1-P)\right)
  =w(M)\frac{M(G-M)}{GP(1-P)}\nabla_zP.
\]
The outer factor $1/G$ and expectation over
$M\sim\operatorname{Binomial}(G,P)$ yield~\eqref{eq:expected-group-gradient}.
For $w\equiv1$,
$\E[M(G-M)]=G(G-1)P(1-P)$, proving the Dr.GRPO value.
\end{proof}

The binomial identity also gives the denominator-free expression
\[
 c_G(P)=\frac{G-1}{G}\sum_{j=0}^{G-2}\binom{G-2}{j}
 w(j+1)P^j(1-P)^{G-2-j}.
\]
The binomial weights sum to one, so $c_G$ lies between
$(G-1)\min_j w(j+1)/G$ and $(G-1)\max_j w(j+1)/G$.
This establishes the smooth, positive extension to $P=0,1$ used in the
subsequent flow and convergence arguments.

The next lemma specializes the practical-estimator analysis of
\citet[Section~4.3 and Appendix~D, Theorem~5]{tajwar2026maxrl}
to our categorical notation; we include the derivation for completeness.

\begin{lemma}[Binary MaxRL has the same correctness-gradient direction]
\label{lem:maxrl-mean}
Let MaxRL assign
$A_i^{\mathrm{MaxRL}}=GR_i/M-1$ when $M>0$ and zero when $M=0$, as in
Section~\ref{sec:maxrl-method}. At the on-policy point, its expected update is
\begin{equation}
  \E[\widehat g_G^{\mathrm{MaxRL}}]
  =c_G^{\mathrm{MaxRL}}(P)\nabla_zP,
  \qquad
  c_G^{\mathrm{MaxRL}}(P)
  =\frac{1-(1-P)^{G-1}}{P}>0,
  \label{eq:maxrl-expected-gradient}
\end{equation}
with continuous endpoint values
$c_G^{\mathrm{MaxRL}}(0)=G-1$ and
$c_G^{\mathrm{MaxRL}}(1)=1$.
\end{lemma}

\begin{proof}
Condition on $M=m>0$. Each successful row has advantage $(G-m)/m$ and each
failed row has advantage $-1$. Using the same conditional score identities as
above, the expected unaveraged sum is
\[
  (G-m)\!\left(\nabla_z\log P-\nabla_z\log(1-P)\right)
  =\frac{G-m}{P(1-P)}\nabla_zP.
\]
After the outer factor $1/G$ and averaging over
$M\sim\operatorname{Binomial}(G,P)$,
\[
 \E[(G-M)\1\{M>0\}]
 =G(1-P)-G\Pr(M=0)
 =G\big[(1-P)-(1-P)^G\big].
\]
Substitution gives Equation~\eqref{eq:maxrl-expected-gradient}. Its endpoint
values follow by continuity (or the finite geometric-series identity), which
also gives the stated continuous extension at both boundaries.
\end{proof}

The same finite-series identity gives the potential used below:
\begin{equation}
 \Psi_{\mathrm{MaxRL}}(P)
 :=\int_0^P c_G^{\mathrm{MaxRL}}(s)\,ds
 =\sum_{k=1}^{G-1}\frac{1-(1-P)^k}{k}.
 \label{eq:maxrl-potential-finite-sum}
\end{equation}
In particular, $\Psi_{\mathrm{MaxRL}}(1)=\sum_{k=1}^{G-1}1/k$.
For binary rewards, best@\(k\) equals pass@\(k\), so this is a harmonic
mixture of sampling objectives for $k=1,\ldots,G-1$. The upper limit is
$G-1$, not $G$: the implemented centered estimator drops the entire update
on all-failure groups. Subtracting an unconditional score baseline even on
those groups instead gives the order-$G$ estimator. This distinction concerns
the exact finite-group expectation at the on-policy point; it does not assert
unbiasedness of subsequent clipped or adaptive-optimizer updates.


\subsection{Winner-take-all collapse in the categorical mean flow}
\label{app:theory-collapse}

Consider $\dot z=c_G(P)\nabla_zP$, with any of the three smooth positive
functions established above. This infinitesimal expected-update model removes
sampling noise and retains independently trained Euclidean logits. The
conditional dynamics reduce to the standard frequency-dependent replicator
system with fitness $f_c(q)=q_c$ \citep{hofbauer1998evolutionary}. The result is
a specialization to the group estimators, not a new general theorem about
replicator dynamics. Related correctness-indifferent collapse analyses include
\citet{sinha2026expected,lochab2026ucpo}.

\begin{theorem}[Standard replicator specialization for binary-reward RL]
\label{thm:grpo-collapse}
Start from finite logits with $0<P(0)<1$. If the conditional correct-mode
distribution $q(0)$ has a unique largest coordinate $c^\star$, then
\[
  P(t)\longrightarrow1,
  \qquad
  q_{c^\star}(t)\longrightarrow1,
  \qquad
  q_c(t)\longrightarrow0\quad(c\ne c^\star).
\]
Thus these GRPO, Dr.GRPO, and binary MaxRL mean flows converge to a single
correct execution mode for generic initial conditions, even though every mode
in $\mathcal C$ has exactly the same reward.
\end{theorem}

\begin{proof}
For a correct mode $c$ and incorrect category $a$,
\[
  \dot z_c=c_G(P)p_c(1-P),
  \qquad
  \dot z_a=-c_G(P)p_aP.
\]
Also $\dot P=c_G(P)\lVert\nabla_zP\rVert_2^2>0$. If $P$ converged to an
interior value, continuity and positivity of $c_G$, together with
\[
 \lVert\nabla_zP\rVert_2^2
 \ge P^2(1-P)^2\left(\frac1m+\frac1n\right),
\]
would keep $\dot P$ bounded away from zero, a contradiction. Hence $P\to1$.

Define the effective optimization time
\[
  \tau(t)=\int_0^t c_G(P(s))P(s)(1-P(s))\,ds.
\]
If $\tau(\infty)<\infty$, every correct and incorrect logit would have finite
total variation, because the absolute velocity of each is at most $\dot\tau$.
All logits would then converge to finite values, contradicting $P\to1$.
Therefore $\tau(t)\to\infty$, leaving infinite effective time for the
correct-simplex dynamics below.

On the correct simplex, changing from $t$ to $\tau$ gives the replicator flow
\begin{equation}
  \frac{dq_c}{d\tau}
  =q_c\left(q_c-\sum_{d\in\mathcal C}q_d^2\right),
  \qquad
  \frac{d}{d\tau}\log\frac{q_c}{q_d}=q_c-q_d.
  \label{eq:grpo-replicator}
\end{equation}
The initially unique maximum remains the unique maximum. For $c\ne c^\star$,
set $\xi_c=q_c/q_{c^\star}$, so $\xi_c(0)<1$ and $\xi_c$ is nonincreasing.
Since $q_{c^\star}\ge1/m$,
\[
 \frac{d\log\xi_c}{d\tau}
 =-q_{c^\star}(1-\xi_c)
 \le-\frac{1-\xi_c(0)}{m}.
\]
Hence $\xi_c(\tau)\le\xi_c(0)\exp[-(1-\xi_c(0))\tau/m]\to0$.
Every competing mode vanishes relative to $c^\star$; normalization then
yields $q_{c^\star}\to1$, establishing the claimed single-mode limit.
\end{proof}

If several modes tie for the largest initial probability, symmetry preserves
the tie; the same ratio argument eliminates strictly smaller coordinates and
leaves the tied maxima uniform. An exactly uniform correct distribution is
therefore a special initial condition. The theorem concerns limiting
concentration: finite softmax logits already give positive probabilities at
every finite time, so finite-time positivity alone is not a retention result.

\paragraph{Dependence on optimization geometry.}
For a smooth parameterization $z=z(\theta)$, Euclidean parameter flow instead
induces $\dot z=c_G(P)J_\theta J_\theta^\top\nabla_zP$. The independently
trained logit theorem uses $J_\theta J_\theta^\top=I$ and does not automatically
extend to arbitrary shared parameters or another optimization metric.
Parameter count alone does not specify this geometry. The empirical scale
comparisons also change pretrained models and learning-rate schedules; they
neither measure the geometry nor identify the cause of cross-scale differences.

\subsection{A replay signal when fresh groups supply none}
\label{app:theory-gradient-availability}

\begin{lemma}[Fresh-group starvation and replay signal]
\label{lem:replay-gradient-availability}
Fix a categorical policy with correctness probability $P$, and draw $G\ge2$
fresh responses independently. For GRPO, Dr.GRPO, and the implemented centered
MaxRL estimator, the probability of a mixed-reward group is
\[
 h_G(P)=1-P^G-(1-P)^G\le G\min\{P,1-P\}.
\]
Every all-equal group has zero fresh task gradient. Thus $h_G$ is the
probability of an opportunity for a nonzero task gradient, rather than a
lower bound on the gradient magnitude.

For a fixed bank with normalized positive weights $w_b$, let
$R_w=-\sum_b w_b\log p_b$ and $\rho_w>0$. Extend $w$ by zero outside the
bank to obtain $\bar w$. The deterministic replay ascent vector is
\[
 v^{\rm replay}=-\rho_w\nabla R_w=\rho_w(\bar w-p).
\]
It vanishes only at $p=\bar w$. In particular, if $p_b\le w_b/2$ for a banked
mode, then $v^{\rm replay}_b\ge\rho_w w_b/2$, even when $b$ is absent from the
fresh group. For a uniform bank of $k\ge2$ modes, as $p\to\mathbf e_b$ at
one banked correct mode, $h_G(P)\to0$ whereas
$\|v^{\rm replay}\|_2\to\rho_w\sqrt{1-1/k}>0$.
The replay norm therefore stays positive as fresh task-signal opportunities vanish.
\end{lemma}

\begin{proof}
The all-correct and all-incorrect events are disjoint and have probabilities
$P^G$ and $(1-P)^G$. A mixed group requires both a correct and an incorrect
sample; the union bound gives the two bounds $GP$ and $G(1-P)$.
For $R_w=-\sum_b w_b\log p_b$, the softmax score identity gives
$\partial R_w/\partial z_a=p_a-\bar w_a$, proving the replay formula and
coordinate bound. Finally,
\[
 \|\bar w-\mathbf e_b\|_2^2=(1-1/k)^2+(k-1)/k^2=1-1/k,
\]
which gives the strictly positive limiting norm for $k\ge2$.
\end{proof}


The displayed replay vector averages the fixed bank. A stochastic exemplar
draw has this vector only in expectation, and a nonzero logit coordinate can
be annihilated by a degenerate neural Jacobian. Replay also need not increase
the next mixed-group probability: $h_G(P)$ is largest at $P=1/2$. The claim
is that a banked target can supply a separate verified signal when the fresh
task gradient is unavailable.

A bounded advantage based on fresh samples has a different boundary behavior.
The next bound includes group-dependent detached scores, and distinguishes
that estimator from exact entropy regularization.

\begin{lemma}[Bounded sampled scores have no probability-independent boundary force]
\label{lem:sampled-score-bound}
Condition on the training history, draw $A_1,\ldots,A_G$ independently from
$p$, and let $a_i$ be any detached, possibly group-dependent advantages with
$|a_i|\le B$. At the on-policy point define
\[
 \widehat g=\frac1G\sum_{i=1}^G a_i(\mathbf e_{A_i}-p).
\]
For every category $b$,
\[
 |\E[\widehat g_b]|\le\E|\widehat g_b|
 \le 2B p_b(1-p_b).
\]
On a group with no draw of $b$, exactly
$\widehat g_b=-p_bG^{-1}\sum_i a_i$; it is zero if the applied advantages
sum to zero. By comparison, the verified-replay component for a fixed target
$w$ has logit velocity $\rho(w_b-p_b)$, which approaches $\rho w_b>0$
for a banked category as $p_b\to0$.
\end{lemma}

\begin{proof}
Use the triangle inequality and $|a_i|\le B$. Each marginal draw satisfies
$\E|\1\{A_i=b\}-p_b|=2p_b(1-p_b)$, irrespective of the dependence of
$a_i$ on the group. On the absence event all indicator terms are zero.
The replay formula follows by differentiating $-\sum_b w_b\log p_b$.
\end{proof}

The bound does not prove inevitable collapse or rule out retention by a
particular bounded estimator. An unsampled mode can still change through
softmax normalization or shared parameters. Exact entropy is outside the
uniformly bounded-score premise because its surprisal is unbounded near zero.

\subsection{Fixed-bank replay and qualitative retention}
\label{app:theory-replay}

Let $\mathcal B\subseteq\mathcal C$ be a nonempty fixed bank of $k$ modes that
the policy has produced and the validator has accepted. Uniform verified replay
uses one exemplar per banked mode and, in the categorical reduction, minimizes
\begin{equation}
  R_{\mathcal B}(z)
  =-\frac1k\sum_{b\in\mathcal B}\log p_b.
  \label{eq:verified-replay-loss}
\end{equation}
This is ordinary cross-entropy against the distribution that assigns mass
$1/k$ to each banked mode and zero elsewhere. It is not conditioned on the
bank: decreasing~\eqref{eq:verified-replay-loss} can raise total bank mass as
well as redistribute mass within the bank.

In an averaged continuous-time model, absorb a fixed round-robin replay
frequency and loss scale into an effective dose $\rho>0$. This assumes a
fixed eligible-bank set and constant relative weighting; it does not identify
finite alternating updates with simultaneous gradient flow. The prompt-isolated categorical mean flow for any of the three fresh
binary-reward objectives above plus verified replay is
\begin{equation}
  \dot z
  =c_G(P)\nabla_zP-\rho\nabla_zR_{\mathcal B}(z).
  \label{eq:replay-mean-flow}
\end{equation}
Let $\Psi(P)=\int_0^P c_G(s)\,ds$. For the finite-group estimators in Lemmas~\ref{lem:grpo-mean}
and~\ref{lem:maxrl-mean}, $c_G$ has a finite positive extension to $[0,1]$;
in particular, Dr.GRPO has $c_G=(G-1)/G$ and MaxRL has
$c_G=[1-(1-P)^{G-1}]/P$. Define
\begin{equation}
  F_{\mathcal B}(z)=\rho R_{\mathcal B}(z)-\Psi(P(z)).
  \label{eq:replay-potential}
\end{equation}
Then~\eqref{eq:replay-mean-flow} is exactly gradient flow
$\dot z=-\nabla_zF_{\mathcal B}$.

\begin{theorem}[Verified replay prevents banked-mode extinction]
\label{thm:replay-retention}
Suppose that after time $T$ the nonempty verified bank $\mathcal B$ is fixed,
fits within replay capacity, and receives a fixed positive effective replay
dose $\rho$. Under~\eqref{eq:replay-mean-flow}, define
\[
  C_T
  =R_{\mathcal B}(z(T))
   +\frac{\Psi(1)-\Psi(P(T))}{\rho}.
\]
Then for every $b\in\mathcal B$ and every $t\ge T$,
\begin{equation}
  p_b(t)\ge \exp(-kC_T)>0.
  \label{eq:replay-retention-bound}
\end{equation}
Thus each banked probability is uniformly bounded away from zero for all
$t\ge T$, including asymptotically, even when the bank does not cover every
correct mode. This is stronger than finite-time positivity, which already
follows from finite softmax logits without replay.
\end{theorem}

\begin{proof}
Along the flow,
\[
  \frac{d}{dt}F_{\mathcal B}(z(t))
  =-\lVert\nabla_zF_{\mathcal B}(z(t))\rVert_2^2\le0.
\]
Since $P(t)\le1$ and $\Psi$ is increasing,
\[
  R_{\mathcal B}(z(t))
  \le R_{\mathcal B}(z(T))
     +\frac{\Psi(1)-\Psi(P(T))}{\rho}
  =C_T.
\]
Every $-\log p_b$ is nonnegative and their sum is
$kR_{\mathcal B}$. Therefore
$-\log p_b(t)\le kC_T$ for every banked mode, which is exactly
Equation~\eqref{eq:replay-retention-bound} for every time after $T$.
\end{proof}


The floor is qualitative and can be numerically vacuous. For example, at
$k=16$ even a hypothetical $C_T=160$ gives $e^{-2560}$, far below any useful
evaluation probability. The argument establishes a uniform-in-time lower
bound in the stated ideal flow, not practical visibility or a prediction for
the experimental replay loss. The effective dose also depends on relative
fresh/replay frequencies, not the replay coefficient alone.

\paragraph{Full coverage and the categorical target.}
\begin{corollary}[Conditional full-coverage categorical limit]
\label{cor:replay-no-collapse}
Assume the conditions of Theorem~\ref{thm:replay-retention} and that by time
$T$ the bank covers the full correct support, $\mathcal B=\mathcal C$. Let
$u_c=1/m$ on $\mathcal C$. Then the categorical mean flow converges to
\[
  P(t)\longrightarrow1,
  \qquad q(t)\longrightarrow u.
\]
Consequently
\[
  \min_{c\in\mathcal C}q_c(t)\longrightarrow\frac1m,
  \qquad
  \mathcal H(q(t))\longrightarrow\log m,
\]
and conditional correct-mode collapse is impossible.
\end{corollary}

\begin{proof}
For full coverage,
\begin{equation}
 R_{\mathcal C}=\log m-\log P+\mathrm{KL}(u\|q).
 \label{eq:replay-mass-uniformity-decomposition}
\end{equation}
The potential is bounded below by $-\Psi(1)$, so its descent gives
$\int_T^\infty\|\nabla F_{\mathcal C}\|_2^2dt<\infty$. The finite-group
$c_G$ is smooth on $[0,1]$, and softmax derivatives and the cross-entropy
gradient are bounded. Thus the Hessian of $F_{\mathcal C}$ and the velocity
are bounded; the gradient along the trajectory is uniformly continuous.
Square integrability then implies $\nabla F_{\mathcal C}\to0$: recurring
gradients of fixed positive size would otherwise contribute energy over
intervals of a uniform minimum length.

Every limiting probability vector therefore satisfies the correct-coordinate
equations
\[
 0=\rho(p_c-1/m)-c_G(P)p_c(1-P).
\]
Summing them gives $0=-(1-P)(\rho+c_G(P)P)$, hence $P=1$. Each correct
coordinate is then $1/m$ and every incorrect coordinate is zero. Compactness
of the probability simplex and uniqueness of this limit establish convergence
of the probability vector.
\end{proof}

\paragraph{Positive weights and the implemented length normalization.}
Uniformity of the target is separate from protection against extinction.
For fixed weights $w_b>0$ with $\sum_{b\in\mathcal B}w_b=1$, define
$R_w=-\sum_b w_b\log p_b$ and use replay coefficient $\rho_w>0$.
The same potential argument gives, for all $t\ge T$,
\[
 p_b(t)\ge\exp(-C_{T,w}/w_b),\qquad
 C_{T,w}=R_w(z(T))+\frac{\Psi(1)-\Psi(P(T))}{\rho_w}.
\]
If $\mathcal B=\mathcal C$, then
$R_w=H(w)-\log P+\mathrm{KL}(w\|q)$, and the independent-logit limit is
$P\to1$, $q\to w$. The preceding stationarity proof applies with $1/m$
replaced by $w_c$: summing the correct-coordinate equations forces $P=1$,
then each equation forces $p_c=w_c$.

A partial bank has a stronger limitation in this same ideal flow. For any
fixed nonempty $\mathcal B\subseteq\mathcal C$ and fixed positive weights
summing to one, extend $w$ by zero to $\bar w$ on all categories. Then
$p(t)\to\bar w$, so even correct modes outside the frozen bank vanish
asymptotically. For $F=\rho_wR_w-\Psi(P)$, the preceding energy and
bounded-Hessian argument still gives $\nabla F\to0$; neither argument
requires full bank coverage.
At every subsequential probability limit, the correct-coordinate equations are
\[
 0=\rho_w(p_c-\bar w_c)-c_G(P)p_c(1-P).
\]
Summing over correct categories gives
$0=-(1-P)(\rho_w+c_G(P)P)$, hence $P=1$; each equation then gives
$p_c=\bar w_c$, and all incorrect mass is zero. Compactness and uniqueness
of this limiting probability vector prove convergence. Thus protecting a
fixed set of banked modes need not preserve unbanked correct alternatives.
This conclusion concerns the fixed-bank categorical flow, not changing-bank
or stochastic neural training.

This distinction matters for the implemented mean-token exemplar loss.
Even if each key were represented by exactly one categorical exemplar of
fixed length $L_b$, its loss would be $A R_w$, where
\[
 A=\frac1k\sum_{b\in\mathcal B}\frac1{L_b},\qquad
 w_b=\frac{L_b^{-1}}{\sum_{j\in\mathcal B}L_j^{-1}}.
\]
Its effective categorical replay coefficient would be $\rho_w=\rho_{\rm LM}A$.
The full-coverage limit would therefore weight shorter exemplars more, and
would be uniform only for equal lengths. For an actual language model,
$\pi_\theta(e_b\mid x)$ is only one contribution to the total key probability.
Equal averaging of per-token exemplar scores does not imply uniform learned
key probabilities; the categorical uniformity corollary is not that claim.

Full support is an additional premise, not a consequence of replay. In
Python factors, prompts admit 16--3,600 correct execution modes (mean 229.44),
while the bank capacity is 16. The full-coverage corollary therefore cannot
supply a domain-wide guarantee there. The retention theorem protects only
admitted banked modes; it supplies no probability floor for an unbanked key.

\subsection{Complete responses and finite evaluation visibility}
\label{app:theory-exemplar-bridge}

The learner scores a particular response string, while evaluation groups
strings by execution key. Event containment provides a conditional bridge:
one complete verified response is one way of producing its key. This bridge
also makes the loss's length factor explicit.

\begin{lemma}[Retention from a joint exemplar potential]
\label{lem:shared-exemplar-retention}
Fix finitely many verified complete-response exemplars $(x_j,e_j)$, each
realizing key $b_j$, with lengths $L_j>0$ and fixed weights $\alpha_j>0$.
Let
\[
 F(\theta)=\sum_j\alpha_j
 \frac{-\log\pi_\theta(e_j\mid x_j)}{L_j}-J(\theta),
 \qquad J(\theta)\le J_{\max}<\infty.
\]
Suppose $F(\theta(T))$ is finite and
$F(\theta(t))\le F(\theta(T))$ for every $t\ge T$. With
$D=F(\theta(T))+J_{\max}$, every protected exemplar satisfies
\[
 p_\theta(b_j\mid x_j)\ge\pi_\theta(e_j\mid x_j)
 \ge\exp\!\left(-\frac{L_jD}{\alpha_j}\right)>0
 \qquad(t\ge T).
\]
\end{lemma}
\begin{proof}
The weighted exemplar loss equals $F(\theta(t))+J(\theta(t))\le D$.
Each summand is nonnegative, so
$-\log\pi_\theta(e_j\mid x_j)\le L_jD/\alpha_j$.
The complete-response event is a subset of the event producing key $b_j$,
which proves the stated lower bound on every protected key's probability.
\end{proof}

Exact descent of a smooth joint potential is sufficient for the energy
inequality, even with shared parameters and multiple prompts. For example,
the ideal categorical model permits
$J(\theta)=\sum_x\nu_x\Psi_x(P_x(\theta))$ with fixed nonnegative prompt
weights and $J_{\max}=\sum_x\nu_x\Psi_x(1)$. The lemma does not establish
this descent condition for stochastic, clipped, alternating AdamW updates,
and it neither implies uniform neural key probabilities nor requires them.

The scored probability must describe the complete response, including its
termination event or a fixed deterministic horizon, under the same prompt,
temperature, masks, truncation rules, and verifier as the sampling law being
bounded. A prefix score or score under another decoding law is insufficient.
For a complete response of length $L$ and mean-token log score $s$, its
probability is $e^{Ls}$, not $e^s$. A change in this exemplar's likelihood
need not equal the change in its whole key's probability.

\paragraph{Retention and finite evaluation visibility.}
At a fixed checkpoint, suppose $k$ distinct banked keys for one prompt have
probability floors $p_b\ge\delta_b>0$. For $K$ independent draws from that
same policy, the banked distinct count $D_{\mathcal B,K}$ satisfies
\[
 \begin{aligned}
 \mathbb E[D_{\mathcal B,K}]
 &\ge\sum_{b\in\mathcal B}\left[1-(1-\delta_b)^K\right],\\
 \Pr(\text{some banked key is unseen})
 &\le\sum_{b\in\mathcal B}(1-\delta_b)^K
 \le k e^{-K\delta_{\min}},
 \end{aligned}
\]
where $\delta_{\min}=\min_b\delta_b$. The first statement follows from
indicator expectations and the second from the union bound. Thus
$K\ge\log(k/\varepsilon)/\delta_{\min}$ suffices to observe all banked keys
with probability at least $1-\varepsilon$, for $0<\varepsilon<1$.
The energy floors can make this budget impractically large; asymptotic
retention does not imply visibility in an eight-draw evaluation.

Conversely, not observing a mode does not establish extinction. At a fixed
checkpoint, zero sightings of a prespecified key in $N$ independent draws
give a one-sided $(1-\alpha)$ binomial upper bound
$1-\alpha^{1/N}$ and lower bound zero \citep{clopper1934binomial}. For
$N=32$ and $\alpha=.05$, the upper bound is approximately $8.94\%$.
These are single-key bounds; simultaneous statements need error allocation.
Counts from changing checkpoints cannot be pooled to certify a current
policy, and teacher-forced scores are not independent Bernoulli observations.

\subsection{Scope relative to entropy, admission, and neural optimization}
\label{app:theory-entropy}
\label{app:theory-entropy-comparison}

\paragraph{Exact entropy is a different comparator.}
For a known full correct support and uniform $u$, the elementary identities
\[
 H(q)=\log m-\mathrm{KL}(q\|u),\qquad
 R_u(p)=\log m-\log P+\mathrm{KL}(u\|q)
\]
show that exact conditional entropy and categorical replay both favor a
uniform conditional target. With strictly increasing correctness potential
$\Psi$, both $\Psi(P)+\beta H(q)$ and $\Psi(P)-\rho R_u(p)$ have the unique
global optimum $P=1,q=u$ on the simplex with $P>0$, for positive coefficients. This follows directly
from nonnegativity of KL and increasing $\Psi$; it does not assert
convergence of either neural update. The KL direction distinguishes their
loss level sets \citep{pereyra2017confidencepenalty,agarwal2021policygradient}:
replay loss diverges when a protected coordinate vanishes, whereas
conditional entropy can remain $\log(m-1)$ with one missing mode.

Exact full-output entropy is a third objective and also regularizes incorrect
mass. Exact tabular entropy-regularized policy gradient has its own positive
probability and convergence guarantees under specified assumptions
\citep{mei2020softmaxpg}; our collapse theorem contains no entropy term.
The sampled semantic scores in the experiments are not exact gradients of
full-output or conditional verified-key entropy. Neither their empirical
comparison nor the bounded-score lemma is an impossibility result for entropy.

\paragraph{Discovery and optimization remain separate requirements.}
Retention starts after a verified key enters active memory. Actual insertion
depends on generation, validation, capacity, and eligibility; a positive
sampling probability alone does not establish eventual admission. Once a
non-evicting capacity-16 bank fills, an unbanked key can remain ineligible.
Changing membership or exemplar weights changes the loss being bounded, so
fixed-bank descent cannot simply be carried across those changes.

For noisy discrete updates, a valid extension would require control of the
joint loss's upward drift, optimizer bias, and accumulated error. A positive
replay coefficient or small nominal learning rate does not establish these
conditions. Current-gradient Adam moments, clipping, momentum, prompt
scheduling, and bank changes are not covered by the exact-flow proofs.
The experiments have not established the required joint-energy or
update-error assumptions. The mathematical contribution is the stated
categorical mechanism and conditional exemplar argument; empirical support
and longitudinal exemplar scores remain separate measurements.
